\section{Referencias Bibliográficas Exhausitvas}

El diseño, arquitectura, desarrollo y documentación de este sistema están respaldados técnica y académicamente por las siguientes directrices oficiales y marcos normativos de ingeniería:

\subsection{Estándares de Ingeniería de Software}
\begin{itemize}
    \item \textbf{IEEE Std 830-1998.} (1998). \textit{IEEE Recommended Practice for Software Requirements Specifications}. Institute of Electrical and Electronics Engineers. (Documento rector utilizado para delinear los requerimientos funcionales y no funcionales del producto y los perfiles de usuario).
    \item \textbf{IEEE Std 1016-2009.} (2009). \textit{IEEE Standard for Information Technology—Systems Design—Software Design Descriptions}. Institute of Electrical and Electronics Engineers. (Empleada para modelar la separación de componentes, diagramas arquitectónicos y la estructura modular de flujos e interfaces de la aplicación).
\end{itemize}

\subsection{Marcos de Trabajo y Herramientas (Frontend)}
\begin{itemize}
    \item \textbf{React Core Team.} (2024). \textit{React: The library for web and native user interfaces. Architecture, Hooks \& Context}. Meta Platforms, Inc. Recuperado de \texttt{https://react.dev/}
    \item \textbf{Evan You, et al.} (2024). \textit{Vite: Next Generation Frontend Tooling. HMR and esbuild integration}. Recuperado de \texttt{https://vitejs.dev/guide/}
    \item \textbf{Tailwind Labs.} (2024). \textit{Tailwind CSS Documentation - Utility-First Fundamentals}. Recuperado de \texttt{https://tailwindcss.com/docs}
    \item \textbf{Tantau, T.} (2023). \textit{TikZ \& PGF Manual for version 3.1.10}. (Extensamente utilizado para la codificación matemática y el renderizado procedural de todos los diagramas de flujo, casos de uso y arquitecturas incrustados en este documento LaTeX).
\end{itemize}

\subsection{APIs e Inteligencia Artificial}
\begin{itemize}
    \item \textbf{Groq API Documentation.} (2024). \textit{Groq Cloud REST API - Fast AI Inference on LPU Architecture}. Groq, Inc. Recuperado de \texttt{https://console.groq.com/docs/quickstart}
    \item \textbf{OpenAI API Reference.} (2024). \textit{Chat Completions Protocol (Adopción de estándares de payload de la industria)}. OpenAI, LLC. Recuperado de \texttt{https://platform.openai.com/docs/api-reference}
    \item \textbf{Crossref.} (2024). \textit{Crossref REST API Documentation for Bibliographic Metadata}. Recuperado de \texttt{https://api.crossref.org/swagger-ui/index.html}
    \item \textbf{Meta AI Research.} (2024). \textit{LLaMA 3: Open Foundation and Fine-Tuned Chat Models}. Meta Platforms, Inc. (Referencia científica sobre el modelo base de 70 Billones de parámetros que dota de inteligencia a los requerimientos funcionales del sistema).
\end{itemize}
